\section{Method}

VLOG-VLA augments a shared flow-matching VLA with a discrete controller state.
We describe the active \texttt{QwenUniversalVLOG} implementation at commit
\texttt{4111898}; graph, critic, learned termination, and legacy Qwen-OFT paths
are not part of this method.

\subsection{Problem Formulation}

At policy query \(t\), an embodiment \(e\) supplies images \(o_t\), instruction
\(\ell\), and normalized native robot state \(q_t^e\).  The policy predicts an
action chunk
\begin{equation}
    A_t^e=a^e_{t:t+H-1}\in\mathbb{R}^{H\times d_a^e}.
    \label{eq:action-chunk}
\end{equation}
The environment may execute only a prefix before the next query.  We distinguish
the decoder horizon \(H\), primitive environment steps, and option age measured
in policy decisions.

The latent controller chooses
\begin{equation}
    z_t\in\{1,\ldots,K\},\qquad
    C=\{c_k\in\mathbb{R}^{d_z}\}_{k=1}^{K},
    \label{eq:option-space}
\end{equation}
and may retain \(z_t\) across several action chunks.  A code is called an
\emph{option} only in this operational sense; no semantic skill label is
assumed.

\begin{figure*}[t]
\centering
\fbox{\parbox{0.95\textwidth}{\centering
\textbf{FIGURE PLACEHOLDER F1: UniversalVLOG architecture}\\[0.45em]
\textbf{Left:} native RoboCasa/LIBERO observations, instruction, state, and
embodiment ID. \quad
\textbf{Center:} shared Qwen3-VL and shared flow-matching DiT, surrounded by
native state/action adapters.\\[0.25em]
\textbf{Upper training-only branch:} encoded future native actions
\(\rightarrow\) posterior \(\rightarrow\) VQ codebook. \quad
\textbf{Inference branch:} current-context router \(\rightarrow\)
per-environment persistent option.\\[0.25em]
\textbf{Action boundary:} explicit option token + bounded residual FiLM on
noisy action tokens \(\rightarrow\) native action decoder.\\[0.25em]
Use gray for pretrained/shared modules, blue for trainable option modules,
purple for native adapters, orange dashed arrows for training-only inputs, and
green for persistent controller state.}}
\caption{\textbf{VLOG-VLA overview.}  One Qwen3-VL backbone and one
flow-matching DiT are shared across embodiments; small native adapters preserve
each robot's state and action coordinates.  Future actions supervise the
posterior only during training.  At inference, a current-context router and
per-environment controller select a persistent code that enters the action
expert through an option token and action-token FiLM.}
\label{fig:architecture}
\end{figure*}


\subsection{Shared Flow-Matching Action Expert}

A Qwen3-VL backbone encodes images and language into
\(H_t^{\mathrm{vl}}\in\mathbb{R}^{L\times d_{\mathrm{vl}}}\).  The action
expert contains one shared DiT, surrounded by native embodiment boundaries:
\begin{align}
    S_t^e &= E_s^e(q_t^e)\in\mathbb{R}^{N_s^e\times d_D},\\
    U_\tau^e &= E_a^e(\widetilde A_\tau^e,\tau)
        \in\mathbb{R}^{H\times d_D},\\
    \widehat V_\tau^e &= D_a^e\!\left(
        \operatorname{DiT}_{\theta}
        (S_t^e,U_\tau^e,H_t^{\mathrm{vl}},\tau)\right).
    \label{eq:native-boundaries}
\end{align}
Here \(E_s^e,E_a^e,D_a^e\) are small native state encoder, action encoder,
and action decoder modules; \(d_D\) is the shared DiT width.  The base
RoboCasa-GR1 embodiment retains the historical GR00T module paths.  Other
embodiments map directly from native normalized coordinates to \(d_D\); for
example, LIBERO's seven-dimensional action is not embedded through the
twenty-nine-dimensional GR1 physical action space.

For a demonstrated action \(A^e\), noise \(\epsilon\), and interpolation time
\(\tau\in[0,1]\), the code constructs
\begin{equation}
    \widetilde A_\tau^e=(1-\tau)\epsilon+\tau A^e,\qquad
    V^{e,*}=A^e-\epsilon,
    \label{eq:flow-path}
\end{equation}
and minimizes a masked flow-matching loss
\begin{equation}
    \mathcal{L}_{\mathrm{FM}}
    =\frac{\sum_{h=1}^{H}m_h
      \|\widehat V_{\tau,h}^e-V_h^{e,*}\|_2^2}
      {\sum_{h=1}^{H}m_h},
    \label{eq:flow-loss}
\end{equation}
where \(m_h\) excludes padding for shorter chunks.  Inference integrates the
predicted velocity for a configured number of Euler steps.

\subsection{Universal Option Discovery}

The controller aggregates the current multimodal representation, native state
tokens, and an embodiment embedding:
\begin{equation}
    s_t=\operatorname{Agg}\!\left(
       P_{\mathrm{vl}}H_t^{\mathrm{vl}},
       \operatorname{mean}(S_t^e)\right)
       +P_e r_e \in\mathbb{R}^{d_z}.
    \label{eq:state-summary}
\end{equation}
During training only, future native actions are passed through the same
embodiment action encoder.  Let
\(F_t^e=E_a^e(A_t^e,\tau{=}0)\) and let \(m_h\) be the action mask.  The
posterior uses four masked temporal summaries:
\begin{equation}
\begin{aligned}
    \mu_t &= \operatorname{mean}_{m} F_t^e,\qquad
    \sigma_t = \operatorname{std}_{m} F_t^e,\\
    \Delta_t &= F_{t,h_{\mathrm{last}}}^e - F_{t,h_{\mathrm{first}}}^e,\\
    \eta_t &= \max_{h:m_hm_{h+1}=1}\left|F_{t,h+1}^e-F_{t,h}^e\right|.
\end{aligned}
\label{eq:temporal-summary}
\end{equation}
The original mean path is preserved for checkpoint compatibility, while a
zero-initialized dynamics projection learns from
\([\sigma_t;\Delta_t;\eta_t]\).  A posterior produces
\begin{equation}
    \begin{aligned}
    y_t&=q_\phi\!\left(s_t,
       \mu_t,\sigma_t,\Delta_t,\eta_t\right),\\
    k_t^+&=\arg\min_k\|y_t-c_k\|_2^2.
    \end{aligned}
    \label{eq:future-posterior}
\end{equation}
The straight-through quantized vector is
\begin{equation}
    \bar c_t=y_t+\operatorname{sg}(c_{k_t^+}-y_t).
    \label{eq:straight-through}
\end{equation}
Codebook and commitment losses train the quantizer.  A motion decoder predicts
the pooled encoded future action from \([s_t;\bar c_t]\).  A differentiable
soft-assignment histogram \(\bar q\) is regularized with an entropy floor,
\begin{equation}
    \mathcal{L}_{\mathrm{use}}
    =
    \left[\rho_H\log K-H(\bar q)\right]_+^2,
    \label{eq:usage-entropy-floor}
\end{equation}
which discourages collapse without forcing every batch to use options
uniformly.  A cosine separation penalty discourages nearly identical code
vectors.

To reduce duration bias during oracle discovery, the implementation can
classify native action chunks into coarse temporal-change strata.  Class zero
denotes low motion; other classes correspond to configured action-dimension
groups, such as arms, hands, and waist for RoboCasa-GR1 or translation,
rotation, and gripper for LIBERO.  These strata are never option labels.  They
produce mixed inverse-frequency sample weights
\begin{equation}
    w_i=(1-\eta_{\mathrm{bal}})+\eta_{\mathrm{bal}}
    \frac{(n_{g_i}+1)^{-\alpha}}
    {\mathbb{E}_j[(n_{g_j}+1)^{-\alpha}]},
    \label{eq:event-balance}
\end{equation}
clipped by a configured maximum.  The weights apply to the VQ, commitment,
motion, and differentiable usage terms in event-balanced oracle stages.

The deployment router sees no future action:
\begin{equation}
    \begin{aligned}
    p_\psi(k\mid s_t)
      &=\operatorname{softmax}(R_\psi(s_t,C))_k,\\
    \mathcal{L}_{R}
      &=\frac{\sum_i \omega_i[-\log p_\psi(k_i^+\mid s_i)]}
      {\sum_i\omega_i}.
    \end{aligned}
    \label{eq:router-loss}
\end{equation}
When router balancing is enabled, \(\omega_i\) uses the same checkpointed
inverse-frequency form over posterior option IDs; otherwise \(\omega_i=1\).
Training and evaluation keep the oracle posterior and the hard deployment
router distinct; a soft mixture is not substituted for the routed option when
measuring deployment action loss.

\subsection{Option-Specific Action Conditioning}

A central design requirement is that changing the code can change the action
expert.  In repair mode, the conditioner first computes
\begin{equation}
    x_k=P_z(\operatorname{LN}(c_k)),
    \label{eq:option-projection}
\end{equation}
using a bias-free option-only projection.  State, embodiment, and learned
projection biases are excluded from this path, preventing a common residual
adapter from passing the oracle-vs-base test while ignoring option identity.
The conditioner produces
\begin{align}
    [\gamma_k;\beta_k]&=W_f x_k,\\
    U_{\tau,k}^{e\prime}
      &=U_\tau^e+\rho\left[
      \tanh(\gamma_k)\odot\operatorname{LN}(U_\tau^e)
      +\tanh(\beta_k)\right],\\
    u_k^{\mathrm{opt}}&=\rho\tanh(W_o x_k),
    \label{eq:option-conditioner}
\end{align}
where the final FiLM and option-token projections are zero initialized and
\(\rho\) is a fixed residual scale.  With fusion disabled, the base embodiment
uses exactly its historical token sequence.  With fusion enabled, the DiT
receives state tokens, an embodiment token for non-base embodiments, the option
token \(u_k^{\mathrm{opt}}\), learned future queries, and conditioned noisy
action tokens.

\begin{figure*}[t]
\centering
\fbox{\parbox{0.95\textwidth}{\centering
\textbf{FIGURE PLACEHOLDER F2: Option-specific conditioning and paired
flow interventions}\\[0.45em]
\textbf{Panel (a):} code \(c_k\) \(\rightarrow\) bias-free option projection
\(\rightarrow\) bounded \(\gamma_k,\beta_k\) and option token; show zero
initialization and fixed \(\rho\).\\[0.25em]
\textbf{Panel (b):} DiT sequence
\([\,\text{state},\text{optional embodiment},\text{option},
\text{future queries},\text{noisy actions}\,]\).\\[0.25em]
\textbf{Panel (c):} one observation/target/mask shared across base, correct,
wrong, shuffled, and hard-router branches with identical \(\tau\) and noise;
label \(\mathcal{L}_0,\mathcal{L}_+,\mathcal{L}_-,
\mathcal{L}_{\mathrm{shuffle}},\mathcal{L}_{R}\).}}
\caption{\textbf{Causal option test at the action-expert boundary.}
Repair-mode conditioning removes state-, embodiment-, and bias-only shortcuts.
Paired branches reuse flow time and noise, so differences isolate the selected
code rather than stochastic flow sampling.}
\label{fig:option-conditioning}
\end{figure*}


The option intervention is paired.  Base, oracle, wrong, shuffled, and hard
router passes reuse the same observation, target, action mask, \(\tau\), and
\(\epsilon\).  Let
\(\mathcal{L}_{0},\mathcal{L}_{+},\mathcal{L}_{-}\) denote fusion-off,
oracle-option, and guaranteed-wrong-option flow losses.  We use
\begin{align}
    \mathcal{L}_{\mathrm{imp}}
      &=\left[m_{\mathrm{imp}}
         +\mathcal{L}_{+}-\mathcal{L}_{0}\right]_+,\\
    \mathcal{L}_{\mathrm{cf}}
      &=\left[m_{\mathrm{cf}}\operatorname{sg}(\mathcal{L}_{0})
         +\mathcal{L}_{+}-\mathcal{L}_{-}\right]_+ .
    \label{eq:paired-losses}
\end{align}
The first term prevents code conditioning from being rewarded by merely
worsening the base; the second requires the correct code to outperform a
different code by a base-relative margin.

\subsection{Persistent Semi-Markov Controller}

Each rollout stream \(i\) stores
\begin{equation}
    M_t^i=(z_t^i,d_t^i,\mathrm{episode}_t^i),
    \label{eq:controller-state}
\end{equation}
where \(d_t^i\) is option age in policy decisions.  A changed episode identifier
clears the active code, age, and cached previous chunk.  Otherwise, the router
candidate \(\tilde z_t=\arg\max_k R_\psi(s_t,C)_k\) is processed as follows:
\begin{equation}
z_t=
\begin{cases}
\tilde z_t, & \text{after reset},\\
z_{t-1}, & d_{t-1}<d_{\min},\\
\tilde z_t, & d_{t-1}\ge d_{\max},\\
\tilde z_t, &
R_{\tilde z_t}-R_{z_{t-1}}>\delta
\ \land\ \tilde z_t\ne z_{t-1},\\
z_{t-1}, & \text{otherwise}.
\end{cases}
\label{eq:controller-rule}
\end{equation}
At \(d_{\max}\) the controller \emph{reselects}; it need not force a different
code.  Age resets only when the active code changes.  This rule implements
cross-query persistence without a learned termination or value model.

\begin{figure}[t]
\centering
\fbox{\parbox{0.92\columnwidth}{\centering
\textbf{FIGURE PLACEHOLDER F3: Persistent controller}\\[0.4em]
Episode reset \(\rightarrow\) initialize router argmax\\
\(\downarrow\)\\
\(d<d_{\min}\): retain active code\\
\(\downarrow\)\\
\(d_{\min}\le d<d_{\max}\): switch only if router-logit margin
\(>\delta\)\\
\(\downarrow\)\\
\(d\ge d_{\max}\): reselect (same code remains possible)\\[0.25em]
Inset: one colored option band spanning several action chunks; mark decision
age and environment-step duration separately.}}
\caption{\textbf{Explicit semi-Markov persistence.}  Independent state is
maintained for each rollout stream and reset by episode ID.  Maximum age causes
reselection, not a forced change, and no learned termination or value signal is
used.}
\label{fig:persistence}
\end{figure}


\subsection{Staged Optimization and Evidence Gates}

The active training schedule separates gradient paths:
\begin{description}
    \item[U0-base:] diagnostic only; all parameters frozen, VLOG disabled.
    \item[U0-native:] train non-base native adapters and the embodiment
    embedding while VLOG remains disabled; the LIBERO configuration names this
    stage \texttt{u0\_libero}.
    \item[U1-oracle:] freeze Qwen, the shared DiT, native adapters, and router;
    train the aggregator, posterior, codebook, motion decoder, and conditioner.
    \item[U2-router:] freeze every module except the router.  Ordinary steps use
    only router cross-entropy; periodic diagnostics evaluate the hard routed
    option through the fixed action expert.
    \item[U3-joint:] keep Qwen frozen and jointly tune the shared DiT, native
    adapters, conditioner, option modules, and router with conservative
    learning rates.
\end{description}

For U1/U3, the implemented objective is
\begin{equation}
\begin{split}
    \mathcal{L}={}&
      \lambda_{\mathrm{FM}}\mathcal{L}_{+}
      +\lambda_{\mathrm{VQ}}\mathcal{L}_{\mathrm{VQ}}
      +\lambda_{\mathrm{com}}\mathcal{L}_{\mathrm{com}}\\
      &+\lambda_{\mathrm{use}}\mathcal{L}_{\mathrm{use}}
      +\lambda_{\mathrm{sep}}\mathcal{L}_{\mathrm{sep}}
      +\lambda_{\mathrm{mot}}\mathcal{L}_{\mathrm{mot}}\\
      &+\lambda_{\mathrm{cf}}\mathcal{L}_{\mathrm{cf}}
      +\lambda_{\mathrm{imp}}\mathcal{L}_{\mathrm{imp}}
      +\mathbb{1}_{\mathrm{U3}}\lambda_R\mathcal{L}_R .
    \label{eq:total-objective}
\end{split}
\end{equation}
Counterfactual frequency is indexed by optimizer step rather than raw forward
calls, and repeated flow samples are disabled in U2 because router-only
optimization has no FM objective.

\begin{figure*}[t]
\centering
\fbox{\parbox{0.95\textwidth}{\centering
\textbf{FIGURE PLACEHOLDER F4: Training curriculum and evidence gates}\\[0.45em]
\begin{tabular}{c|c|c}
\textbf{Stage} & \textbf{Trainable path} & \textbf{Required evidence}\\
\hline
U0-base & none & strict load, fusion-off parity, base rollout\\
U0-native & native adapters + embodiment embedding & native rollout, base retention\\
U1-oracle & aggregator/posterior/codebook/conditioner & Gate O: base/correct/wrong/shuffle\\
U2-router & router only & Gate R: hard-router recovery/action loss\\
U3-joint & shared DiT/adapters/VLOG; Qwen frozen & paired rollout non-inferiority/benefit
\end{tabular}\\[0.4em]
Below the stages, show immutable artifacts: resolved config, checkpoint, gate
JSON, per-episode result, option timeline, and profiler report.}}
\caption{\textbf{Separated optimization and promotion gates.}  The schedule
prevents router training from changing an oracle mapping that has already
passed Gate O.  Offline flow gates validate mechanisms; simulator rollouts are
still required for task claims.}
\label{fig:objectives}
\end{figure*}


U1 advances only if Gate O shows a positive correct-vs-wrong and
correct-vs-shuffled gap, non-collapsed use, a nonzero but bounded residual, and
no degradation relative to the fusion-off base.  U2 advances only if Gate R
shows that the hard router recovers active codes and remains close to oracle
action loss.  These are paired offline mechanism gates, not rollout results.

\subsection{Inference}

At deployment, future actions and the posterior are absent.  Qwen3-VL and the
native state adapter produce current features, the hard router proposes a code,
the per-environment controller updates its persistent state, and the selected
code conditions every Euler velocity evaluation of the shared flow expert.
The native decoder returns a normalized action chunk in embodiment coordinates.
The deployment wrapper must supply stable environment and episode identifiers
so controller state cannot leak across rollout streams or episodes.
